\chapter[Lezione 9 --- 16 ottobre 2026]{Lezione 9: limiti all'infinito, composizione e applicazioni inverse}
\label{chap:lezione-09}

\section[Verifica di un limite all'infinito]{Verifica di un limite per \(x\to+\infty\)}
\label{sec:lez09-limite-infinito}

Si completa lo studio degli esempi della \cref{sec:lez08-esercizi-limiti}
considerando un limite finito per \(x\to+\infty\).

\begin{exercise}[Verifica di un limite all'infinito]
\label{ex:lez09-limite-razionale-infinito}
Verificare mediante la \cref{def:lez08-limite} che
\[
  \lim_{x\to+\infty}\frac{3x-1}{x+2}=3.
\]
\end{exercise}

\begin{proof}[Soluzione dell'esercizio]
\emph{Dominio e punto di accumulazione.}
La funzione
\[
  f(x):=\frac{3x-1}{x+2}
\]
ha dominio
\[
  A=\R\setminus\{-2\}
   =(-\infty,-2)\cup(-2,+\infty).
\]
Poiché \(A\) non è superiormente limitato,
\(+\infty\) è un punto di accumulazione per \(A\)
in virtù della \cref{prop:lez08-accumulazione-limitatezza}.

\emph{Verifica della definizione.}
Fissato \(\varepsilon>0\), occorre trovare
\(k_\varepsilon\in\R\) tale che
\[
  \left|\frac{3x-1}{x+2}-3\right|<\varepsilon
  \qquad\forall x\in A\cap(k_\varepsilon,+\infty).
\]
Per ogni \(x\in A\), si ha
\[
  \left|\frac{3x-1}{x+2}-3\right|
  =\left|\frac{-7}{x+2}\right|
  =\frac7{|x+2|}.
\]
Se \(x>-2\), allora \(|x+2|=x+2\), e la condizione
\[
  \frac7{x+2}<\varepsilon
\]
è equivalente a \(x>-2+7/\varepsilon\). È pertanto
sufficiente scegliere
\[
  k_\varepsilon:=-2+\frac7\varepsilon.
\]
Infatti, per ogni \(x\in A\cap(k_\varepsilon,+\infty)\)
si ha \(x>-2\) e
\[
  \left|f(x)-3\right|
  =\frac7{x+2}<\varepsilon.
\]
Poiché \(\varepsilon>0\) era arbitrario, per la
\cref{def:lez08-limite} si conclude
\[
  \lim_{x\to+\infty}\frac{3x-1}{x+2}=3.
\]
\end{proof}

\clearpage

\begin{figure}[htbp]
  \centering
  \begin{tikzpicture}[unipd/diagram,scale=0.95]
  \draw[unipd/directed edge] (-4.5,0) -- (6.2,0);
  \node[below] at (6.1,-.02) {\(x\)};
  \coordinate (excluded) at (-2.6,0);
  \coordinate (threshold) at (.65,0);
  \draw[unipd-accent,line width=2.7pt,opacity=.45] (threshold) -- (5.85,0);
  \draw[unipd-accent,line width=.75pt,dashed] (.65,0) -- (.65,1.05);
  \draw[unipd-primary,line width=.75pt] (-2.6,-.10) -- (-2.6,.10);
  \draw[unipd-primary,line width=.75pt] (.65,-.10) -- (.65,.10);
  \fill[white] (-2.6,0) circle (2.9pt);
  \draw[unipd-primary,line width=.75pt] (-2.6,0) circle (2.9pt);
  \fill[white] (.65,0) circle (2.9pt);
  \draw[unipd-accent,line width=.75pt] (.65,0) circle (2.9pt);
  \node[below=6pt] at (-2.6,0) {\(-2\)};
  \node[below=6pt] at (.65,0) {\(k_\varepsilon=-2+\dfrac7\varepsilon\)};
  \node[above,text=unipd-primary] at (-2.6,.18) {punto escluso};
  \node[above,text=unipd-accent!80!black] at (3.15,.18) {\(x>k_\varepsilon\)};
\end{tikzpicture}

  \caption{La soglia \(k_\varepsilon=-2+7/\varepsilon\)
  individua un intorno di \(+\infty\) nel quale
  \(x+2>0\) e \(|f(x)-3|<\varepsilon\).}
  \label{fig:lez09-limite-razionale-infinito}
\end{figure}

\section{Composizione di applicazioni}
\label{sec:lez09-composizione}

\begin{definition}[Applicazione composta]
\label{def:lez09-composizione}
Siano \(A,B,C,D\) insiemi non vuoti e siano
\[
  \varphi\colon A\longrightarrow B,
  \qquad \psi\colon C\longrightarrow D,
\]
tali che \(\operatorname{Im}(\varphi)\subseteq C\).
Si definisce \emph{applicazione composta} di \(\varphi\)
e \(\psi\) l'applicazione
\[
  \psi\circ\varphi\colon A\longrightarrow D,
  \qquad a\longmapsto\psi\bigl(\varphi(a)\bigr).
\]
\end{definition}

L'ordine della composizione è indicato nella
\cref{fig:lez09-composizione}: si applica prima \(\varphi\)
e successivamente \(\psi\).

\begin{figure}[htbp]
  \centering
  \begin{tikzpicture}[unipd/diagram]
  \foreach \pos/\set in {0/A,4.3/C,8.6/D} {
    \draw[unipd/edge] (\pos,0) ellipse (.83 and 1.0);
    \node at (\pos,1.26) {\(\set\)};
    \fill[unipd-primary] (\pos,0) circle (2pt);
  }
  \node[below=7pt] at (0,0) {\(a\)};
  \node[below=7pt] at (4.3,0) {\(\varphi(a)\)};
  \node[below=7pt] at (8.6,0) {\(\psi(\varphi(a))\)};
  \draw[unipd/directed edge] (.25,.20) -- (4.05,.20)
    node[midway,above=3pt] {\(\varphi\)};
  \draw[unipd/directed edge] (4.55,.20) -- (8.35,.20)
    node[midway,above=3pt] {\(\psi\)};
  \draw[unipd-accent,dashed,-{Latex[length=2.2mm]},line width=.85pt]
    (.08,1.46) to[bend left=22]
    node[midway,above=1pt] {\(\psi\circ\varphi\)} (8.52,1.46);
  \node[below=7pt,text=unipd-primary] at (4.3,-1.06)
    {\(\operatorname{Im}(\varphi)\subseteq C\)};
\end{tikzpicture}

  \caption{Composizione \(\psi\circ\varphi\), con
  \(\operatorname{Im}(\varphi)\subseteq D(\psi)=C\).}
  \label{fig:lez09-composizione}
\end{figure}

\begin{example}[L'ordine della composizione]
\label{ex:lez09-ordine-composizione}
Si considerino le applicazioni
\[
  \varphi\colon\R\longrightarrow\R,
  \quad x\longmapsto x^2,
  \qquad
  \psi\colon\R\longrightarrow\R,
  \quad x\longmapsto x+1.
\]
Entrambe le composizioni sono definite e, per ogni \(x\in\R\),
\[
  (\psi\circ\varphi)(x)=\psi(x^2)=x^2+1,
  \qquad
  (\varphi\circ\psi)(x)=\varphi(x+1)=(x+1)^2.
\]
Per \(x=10\), in particolare,
\[
  (\psi\circ\varphi)(10)=101,
  \qquad (\varphi\circ\psi)(10)=121.
\]
Pertanto \(\psi\circ\varphi\ne\varphi\circ\psi\):
la composizione non è, in generale, commutativa.
\end{example}

L'esistenza delle due composizioni richiede inoltre
condizioni distinte, come mostra l'esempio seguente.

\clearpage

\begin{example}[Composizione non definita]
\label{ex:lez09-composizione-non-definita}
Siano
\[
  \varphi\colon\N\longrightarrow\Q,
  \quad n\longmapsto\frac n2,
  \qquad
  \psi\colon\N\longrightarrow\R,
  \quad n\longmapsto n^2.
\]
Poiché \(\operatorname{Im}(\psi)\subseteq\N=D(\varphi)\),
la composizione \(\varphi\circ\psi\) è definita e
\[
  (\varphi\circ\psi)(n)=\frac{n^2}{2}
  \qquad\forall n\in\N.
\]
Viceversa, \(\varphi(1)=1/2\notin\N=D(\psi)\);
pertanto \(\operatorname{Im}(\varphi)\not\subseteq D(\psi)\)
e \(\psi\circ\varphi\) non è definita su tutto \(\N\).
\end{example}

\section{Applicazione identità}
\label{sec:lez09-identita}

\begin{definition}[Applicazione identità]
\label{def:lez09-identita}
Sia \(A\) un insieme non vuoto. Si chiama
\emph{applicazione identità} su \(A\) l'applicazione
\[
  \operatorname{id}_A\colon A\longrightarrow A,
  \qquad a\longmapsto a.
\]
\end{definition}

Per \(A=\R\), il grafico definito nella
\cref{def:lez07-grafico} è
\[
  G(\operatorname{id}_{\R})
  =\{(a,a):a\in\R\}.
\]

\begin{proposition}[Identità e composizione]
\label{prop:lez09-identita-composizione}
Sia \(\varphi\colon A\longrightarrow B\) un'applicazione
tra insiemi non vuoti. Allora
\[
  \varphi\circ\operatorname{id}_A=\varphi,
  \qquad
  \operatorname{id}_B\circ\varphi=\varphi.
\]
\end{proposition}

\begin{proof}
Per ogni \(a\in A\), applicando le
\cref{def:lez09-composizione,def:lez09-identita}, si ha
\[
  (\varphi\circ\operatorname{id}_A)(a)
  =\varphi(\operatorname{id}_A(a))=\varphi(a),
\]
\[
  (\operatorname{id}_B\circ\varphi)(a)
  =\operatorname{id}_B(\varphi(a))=\varphi(a).
\]
Le due composizioni hanno dominio \(A\) e codominio \(B\),
come \(\varphi\), dunque valgono entrambe le identità.
\end{proof}

Considerando la stessa legge di \(\varphi\) con codominio
\(\operatorname{Im}(\varphi)\), si ottiene l'applicazione
\(\widetilde{\varphi}\colon A\to\operatorname{Im}(\varphi)\).
Per la \cref{prop:lez09-identita-composizione} vale anche
\[
  \operatorname{id}_{\operatorname{Im}(\varphi)}
  \circ\widetilde{\varphi}=\widetilde{\varphi}.
\]

\begin{figure}[H]
  \centering
  \begin{tikzpicture}[unipd/diagram]
  \foreach \yy in {0,-2.05} {
    \foreach \xx in {0,3.35,6.7} {
      \draw[unipd/edge] (\xx,\yy) ellipse (.65 and .65);
      \fill[unipd-primary] (\xx,\yy) circle (2pt);
    }
  }
  \node at (0,.9) {\(A\)};
  \node at (3.35,.9) {\(A\)};
  \node at (6.7,.9) {\(B\)};
  \node at (0,-1.15) {\(A\)};
  \node at (3.35,-1.15) {\(B\)};
  \node at (6.7,-1.15) {\(B\)};
  \draw[unipd/directed edge] (.15,.12) -- (3.2,.12)
    node[midway,above] {\(\operatorname{id}_A\)};
  \draw[unipd/directed edge] (3.5,.12) -- (6.55,.12)
    node[midway,above] {\(\varphi\)};
  \draw[unipd/directed edge] (.15,-1.93) -- (3.2,-1.93)
    node[midway,above] {\(\varphi\)};
  \draw[unipd/directed edge] (3.5,-1.93) -- (6.55,-1.93)
    node[midway,above] {\(\operatorname{id}_B\)};
  \foreach \xx in {0,3.35} {
    \node[below=5pt] at (\xx,0) {\(a\)};
  }
  \node[below=5pt] at (6.7,0) {\(\varphi(a)\)};
  \node[below=5pt] at (0,-2.05) {\(a\)};
  \foreach \xx in {3.35,6.7} {
    \node[below=5pt] at (\xx,-2.05) {\(\varphi(a)\)};
  }
\end{tikzpicture}

  \caption{L'applicazione identità è elemento neutro
  a destra e a sinistra per la composizione.}
  \label{fig:lez09-identita-composizione}
\end{figure}

\clearpage

\section{Applicazione inversa}
\label{sec:lez09-inversa}

Per invertire un'applicazione occorre poter risalire
in modo univoco da ogni elemento dell'immagine al suo
antecedente. La \cref{def:lez06-iniettiva} assicura tale unicità.

\begin{example}[Ostacolo all'esistenza dell'inversa]
\label{ex:lez09-non-iniettiva-inversa}
Per l'applicazione
\(\varphi\colon\R\to\R\), \(\varphi(x)=x^2\), si ha
\[
  \varphi(-1)=\varphi(1)=1.
\]
La \cref{def:lez06-controimmagine} fornisce
\[
  \varphi^{-1}(\{1\})=\{-1,1\}.
\]
Il valore \(1\) possiede quindi due antecedenti distinti:
non è possibile definirne un'inversa che associ a \(1\)
un unico elemento di \(\R\).
\end{example}

\begin{figure}[htbp]
  \centering
  \begin{tikzpicture}[unipd/diagram]
  \draw[unipd/edge] (0,0) ellipse (1 and 1.35);
  \draw[unipd/edge] (4.5,0) ellipse (.9 and 1.35);
  \node at (0,1.65) {\(\R\)};
  \node at (4.5,1.65) {\(\R\)};
  \fill[unipd-primary] (0,.65) circle (2pt);
  \fill[unipd-primary] (0,-.65) circle (2pt);
  \fill[unipd-primary] (4.5,0) circle (2pt);
  \node[left=5pt] at (0,.65) {\(-1\)};
  \node[left=5pt] at (0,-.65) {\(1\)};
  \node[right=5pt] at (4.5,0) {\(1\)};
  \draw[unipd/directed edge] (.15,.62) -- (4.35,.04);
  \draw[unipd/directed edge] (.15,-.62) -- (4.35,-.04);
  \node at (2.25,1.05) {\(\varphi(x)=x^2\)};
  \node[align=center,text=unipd-accent!80!black] at (2.25,-1.65)
    {due preimmagini dello stesso valore};
\end{tikzpicture}

  \caption{La funzione \(\varphi(x)=x^2\) non è iniettiva:
  \(-1\) e \(1\) hanno la stessa immagine.}
  \label{fig:lez09-inversa-non-iniettiva}
\end{figure}

\begin{definition}[Applicazione inversa]
\label{def:lez09-inversa}
Siano \(A,B\) insiemi non vuoti e sia
\(\varphi\colon A\longrightarrow B\) iniettiva.
Si chiama \emph{applicazione inversa} di \(\varphi\)
l'applicazione
\[
  \varphi^{-1}\colon\operatorname{Im}(\varphi)
  \longrightarrow A,
  \qquad b\longmapsto a,
\]
dove \(a\in A\) è l'unico elemento tale che
\(\varphi(a)=b\).
\end{definition}

L'esistenza di \(a\) segue dalla
\cref{def:lez06-immagine}, mentre l'unicità segue dalla
\cref{def:lez06-iniettiva}. In particolare,
\(D(\varphi^{-1})=\operatorname{Im}(\varphi)\) e
\(C(\varphi^{-1})=A\).

\begin{figure}[H]
  \centering
  \begin{tikzpicture}[unipd/diagram]
  \foreach \xx/\set in {0/A,4.1/{\operatorname{Im}(\varphi)},8.2/A} {
    \draw[unipd/edge] (\xx,0) ellipse (.75 and 1);
    \node at (\xx,1.32) {\(\set\)};
    \fill[unipd-primary] (\xx,0) circle (2pt);
  }
  \node[below=7pt] at (0,0) {\(a\)};
  \node[below=7pt] at (4.1,0) {\(\varphi(a)\)};
  \node[below=7pt] at (8.2,0) {\(a\)};
  \draw[unipd/directed edge] (.2,.19) -- (3.9,.19)
    node[midway,above=3pt] {\(\varphi\)};
  \draw[unipd/directed edge] (4.3,.19) -- (8,.19)
    node[midway,above=3pt] {\(\varphi^{-1}\)};
  \draw[unipd-accent,dashed,-{Latex[length=2.2mm]},line width=.85pt]
    (.08,-1.2) to[bend right=20]
    node[midway,below] {\(\varphi^{-1}\circ\varphi=\operatorname{id}_A\)}
    (8.12,-1.2);
\end{tikzpicture}

  \caption{L'applicazione inversa riporta
  \(\varphi(a)\) ad \(a\), realizzando
  \(\varphi^{-1}\circ\varphi=\operatorname{id}_A\).}
  \label{fig:lez09-composizione-inversa}
\end{figure}

\begin{proposition}[Composizione con l'applicazione inversa]
\label{prop:lez09-inversa-composizione}
Sia \(\varphi\colon A\longrightarrow B\) iniettiva.
Con la notazione introdotta nella \cref{sec:lez09-identita},
si ha
\[
  \varphi^{-1}\circ\varphi=\operatorname{id}_A,
  \qquad
  \widetilde{\varphi}\circ\varphi^{-1}
  =\operatorname{id}_{\operatorname{Im}(\varphi)}.
\]
\end{proposition}

\begin{proof}
Le due composizioni sono definite rispettivamente
da \(A\) in \(A\) e da \(\operatorname{Im}(\varphi)\)
in sé.

Per ogni \(a\in A\), dalla \cref{def:lez09-inversa}
si ottiene
\[
  (\varphi^{-1}\circ\varphi)(a)
  =\varphi^{-1}(\varphi(a))=a
  =\operatorname{id}_A(a).
\]
Analogamente, per ogni \(b\in\operatorname{Im}(\varphi)\),
\[
  (\widetilde{\varphi}\circ\varphi^{-1})(b)
  =\widetilde{\varphi}(\varphi^{-1}(b))=b
  =\operatorname{id}_{\operatorname{Im}(\varphi)}(b).
\]
Seguono entrambe le uguaglianze.
\end{proof}

\begin{example}[Esponenziale e logaritmo]
\label{ex:lez09-inversa-esponenziale}
Si considerino le applicazioni inverse
\[
  \varphi\colon\R\longrightarrow(0,+\infty),
  \quad x\longmapsto e^x,
\]
\[
  \varphi^{-1}\colon(0,+\infty)\longrightarrow\R,
  \quad y\longmapsto\log y.
\]
Per ogni \(x\in\R\), si ha
\[
  (\varphi^{-1}\circ\varphi)(x)
  =\log(e^x)=x.
\]
Analogamente, per ogni \(y\in(0,+\infty)\),
\[
  (\varphi\circ\varphi^{-1})(y)
  =e^{\log y}=y.
\]
Pertanto
\[
  \varphi^{-1}\circ\varphi=\operatorname{id}_{\R},
  \qquad
  \varphi\circ\varphi^{-1}
  =\operatorname{id}_{(0,+\infty)}.
\]
Le due composizioni hanno domini differenti:
\[
  D(\varphi^{-1}\circ\varphi)=\R,
  \qquad
  D(\varphi\circ\varphi^{-1})=(0,+\infty).
\]
\end{example}
